% ============================================================================
% Part 13: Concrete readability improvement plan (source-grounded)
% ============================================================================
\chapter{Making the generated contract human-readable: a concrete plan}
\label{ch:readability}

Chapters~\ref{ch:prod-review} and~\ref{ch:critique} established that the
artifact is \emph{correct} but hard for a human to audit. This chapter turns
that diagnosis into a concrete, prioritized change plan. Every fix below is
anchored to a specific line in the generator
(\texttt{proofs/\allowbreak{}solidity-verifier/\allowbreak{}src/\allowbreak{}}, \texttt{main} HEAD \texttt{6b3d2096})
and shows the
\emph{before} (what the renderer emits today) and the \emph{after} (what it
should emit). Each fix carries a status tag as in Chapter~\ref{ch:critique}
(\vstatus{ADDRESSED} / \vstatus{PARTIAL} / \vstatus{OPEN} /
\vstatus{CORRECTED}); \texttt{L}$n$ refers to
\texttt{Halo2VerifierFee12345B64.sol}. The central observation that makes
this tractable:

\begin{quote}
\textbf{The generator already knows the semantic name of almost every
address and value it emits --- it just throws the name away at the emit
boundary.}
\end{quote}

\paragraph{Where the names are, and where the hex comes from}
\vstatus{CORRECTED}. The names exist in the data model:
\texttt{MemoryRegion.name} (\texttt{layout/\allowbreak{}memory.rs:179}) and the
\texttt{fixed\_\allowbreak{}low\_\allowbreak{}memory} table (\texttt{memory.rs:943}) are used only for
validation error strings, while \texttt{column\_eval\_var}
(\texttt{yul\_emit.rs:1142}) is \emph{already} used at emit and produces the
\texttt{f\_N}/\texttt{a\_N}/\texttt{i\_N} names. Contrary to an earlier
draft, \texttt{u256\_string} (\texttt{yul\_emit.rs:1133}) renders
\emph{field constants}, not addresses. Addresses reach the Yul through the
\texttt{Ptr}/\texttt{Value} \texttt{Display} impl
(\texttt{encoding/\allowbreak{}mod.rs:634--648}) and through
\texttt{format!("\{:\#x\}")} sites (64 in \texttt{kzg/mod.rs}, 31 in
\texttt{quotient.rs}). \texttt{Value} already has a symbolic variant,
\texttt{Value::Identifier(name, off)} (\texttt{encoding/\allowbreak{}mod.rs:590}), which
renders as \texttt{add(name, off)}; that is how
\texttt{add(X1\_\allowbreak{}POWERS\_\allowbreak{}MPTR, 0x80)} appears in the output. Recovering
readability is therefore mostly a matter of \emph{propagating names that
already exist} to the emit site, not inventing new information.

\section{The measured readability debt}
\label{sec:read-measure}

Before proposing fixes, the debt in the production
\texttt{Halo2VerifierFee12345B64.sol} (3{,}579 lines; \texttt{wc -l} prints
3{,}578). Address metrics use the definition of
\S\ref{sec:c-scorecard}: the first argument of
\texttt{mload}/\texttt{mstore}/\texttt{mcopy}, comments stripped.

\begin{center}
\begin{tabular}{@{}lrl@{}}
\toprule
\textbf{Metric} & \textbf{Count} & \textbf{Verdict} \\
\midrule
Raw-hex address operands (distinct addresses) & 492 (246) & unreadable \\
Named pointer operands, bare / as \texttt{add(NAME, off)} & 64 / 112 & good, but minority \\
Declared pointer constants (referenced) & 68 (63) & 5 are dead \\
Opaque \texttt{varN} local names (uses) & 163 (980) & unreadable \\
Lines in \texttt{verifyProof} (L293--L3578) & 3{,}286 & unreviewable as one unit \\
Region names defined in \texttt{layout/memory.rs} & 30 & \emph{one per buffer, unused at emit} \\
\bottomrule
\end{tabular}
\end{center}

The ratio is stark: for every bare named address operand the reader gets,
they get $\sim$7.7 raw hex ones ($\sim$2.9 if \texttt{add(NAME, off)} forms
count as named). The raw operands cluster in two buffers: 327 operands (183
addresses) in the final-MSM/quotient scratch from \texttt{0xa320}, and 160
(60) in the decoded-eval buffer from \texttt{REVERSED\_\allowbreak{}EVALS\_\allowbreak{}MPTR}
$=$ \texttt{0x8520}; only 5 (3) are low scratch. The fix is to invert that
ratio.

\section{Fix 1 --- Named memory addresses (highest impact)}
\label{sec:read-named-addr}
\vstatus{OPEN} \emph{(prototyped on the unmerged branch \texttt{41b78eb9},
2026-08-20, for the final-MSM block only; not in \texttt{main} or B64)}

\paragraph{The problem.} A typical emitted block (L3455--3458):

\begin{lstlisting}[style=yul]
mcopy(0xd2a0, 0x9c60, 0x80)
mstore(0xd320, mulmod(mload(add(X1_POWERS_MPTR, 0x80)), x4_pow_4, r))
mcopy(0xd340, F_COM_MPTR, 0x80)
mstore(0xd3c0, x4_pow_5)
\end{lstlisting}

The reader can see \texttt{X1\_POWERS\_MPTR} and \texttt{F\_COM\_MPTR}
are named, but \texttt{0xd2a0}, \texttt{0x9c60}, \texttt{0xd320},
\texttt{0xd340}, \texttt{0xd3c0} are opaque. They are pairs 76 and 77 of
the 78-pair final MSM staged at \texttt{0xa320} (point at
$\texttt{0xa320}+76\cdot\texttt{0xa0}$, scalar $+\texttt{0x80}$), and
\texttt{0x9c60} is \texttt{PERM\_\allowbreak{}Z\_\allowbreak{}COMMS\_\allowbreak{}MPTR\_\allowbreak{}BASE}${}+\texttt{0x200}$,
the fifth permutation-product commitment --- but nothing says so.

\paragraph{What is half-built, and what is not.} \texttt{MemoryRegion}
(\texttt{memory.rs:178}) carries a \texttt{name: \&'static str} field,
and \texttt{MemoryMap} holds a \texttt{Vec<MemoryRegion>} (L217--218). Its
doc comment scopes the name to errors: \emph{``Stable human-readable region
name used in validation errors''}. An earlier draft proposed a reverse lookup
\texttt{name\_\allowbreak{}for\_\allowbreak{}addr(addr)} matching \texttt{r.start == addr}. That does
not work, for three reasons:
\begin{itemize}[leftmargin=1.4em,itemsep=1pt]
\item Regions are one per \emph{buffer} (\texttt{"pcs\_final\_msm"},
  \texttt{"decoded\_evals"}, \ldots), so an exact-start match names only the
  base. Every interior slot above --- all five addresses --- still falls
  back to hex. No per-slot \texttt{MSM\_IN\_$i$} regions exist.
\item Region names are lowercase planner labels, not declared Solidity
  identifiers, so emitting one would not compile without a matching constant.
\item Addresses are phase-shared (\texttt{0x80} is \texttt{TRANSCRIPT\_MPTR},
  \texttt{RETURN\_MPTR} and \texttt{QUOTIENT\_\allowbreak{}RETURN\_\allowbreak{}MPTR};
  \texttt{0xa1e0} is \texttt{SELECTOR\_\allowbreak{}ACC\_\allowbreak{}MPTR} and
  \texttt{BATCH\_\allowbreak{}INV\_\allowbreak{}SCRATCH\_\allowbreak{}MPTR}; \texttt{0xa320} is both the final-MSM
  base and a quotient accumulator), so an address alone cannot pick the
  right name. The \emph{emit site} knows the role; the reverse map does not.
\end{itemize}

\paragraph{Concrete change.} Name the pointer where it is created, using the
symbolic pointer the encoder already supports, and declare the base once:

\begin{lstlisting}[style=rust]
// kzg/mod.rs (sketch): stage final-MSM pairs relative to a named base.
// Value::Identifier(name, off) already renders as add(name, off)
// (encoding/mod.rs:634); Constants.sol declares PCS_FINAL_MSM_MPTR.
fn final_msm_slot(pair: usize, scalar: bool) -> Value {
    let off = pair * 0xa0 + if scalar { 0x80 } else { 0 };
    Value::Identifier("PCS_FINAL_MSM_MPTR", off as isize)
}
\end{lstlisting}

Do the same for the decoded-eval reads (\texttt{REVERSED\_\allowbreak{}EVALS\_\allowbreak{}MPTR}
${}+k\cdot\texttt{0x20}$). A containment lookup
(\texttt{r.start <= addr \&\& addr < r.start + r.len}) is still useful, but
only for the generated legend of Fix~5, where ambiguity can be shown rather
than resolved.

\paragraph{After.} The same block becomes:

\begin{lstlisting}[style=yul]
mcopy(add(PCS_FINAL_MSM_MPTR, 0x2f80), add(PERM_Z_COMMS_MPTR_BASE, 0x0200), 0x80) // pair 76: perm_z[4]
mstore(add(PCS_FINAL_MSM_MPTR, 0x3000), mulmod(mload(add(X1_POWERS_MPTR, 0x80)), x4_pow_4, r))
mcopy(add(PCS_FINAL_MSM_MPTR, 0x3020), F_COM_MPTR, 0x80)                          // pair 77: f_com
mstore(add(PCS_FINAL_MSM_MPTR, 0x30a0), x4_pow_5)
\end{lstlisting}

Now the reader sees ``this is the final-MSM staging buffer, pair $i$''
without a memory-map lookup. Converting every emitter that addresses the
\texttt{0xa320} scratch (final-MSM staging in \texttt{kzg/mod.rs},
quotient-VM state in \texttt{quotient.rs}) and the decoded-eval reads removes
487 of the 492 raw-hex operands; the remaining handful are low scratch.

\paragraph{Cost / risk.} Low. The emitted \emph{semantics} are unchanged
(a \texttt{constant} is substituted at compile time and the via-IR optimizer
folds \texttt{add(const, const)}). Byte-identical bytecode is expected but
should be \emph{checked}, not asserted: the crate's release check already
pins the IVC verifier's runtime keccak
(\texttt{scripts/\allowbreak{}check\_\allowbreak{}release\_\allowbreak{}bytecode\_\allowbreak{}sizes.sh}), which turns ``only
the source text changed'' into a test, and the trace-equivalence harness
(\S\ref{sec:c2-trustroot}) must still pass unchanged. Neither is gated by
the monorepo's root CI today, so run them explicitly.

\paragraph{Evidence the mechanism already works.} The contract already
declares 78 \texttt{uint256 internal constant} names (79 with the
\texttt{bytes32} codehash; e.g.\
\texttt{PROOF\_\allowbreak{}CPTR = 0x64}, \texttt{INSTANCE\_\allowbreak{}CPTR = 0x1ee4},
\texttt{NUM\_\allowbreak{}INSTANCE\_\allowbreak{}CPTR = 0x1ec4}) that the Yul body
references by name and that compile to the same literals. An earlier draft
said the MSM scratch base \texttt{0xa320} was the \emph{only} kind of address
still emitted raw; that is wrong --- the decoded-eval buffer accounts for a
third of the raw operands (above). Extending the existing constant mechanism
to the scratch regions is not a new capability --- it is using one the
generator already has. A named \texttt{PCS\_\allowbreak{}FINAL\_\allowbreak{}MSM\_\allowbreak{}MPTR} resolves to
\texttt{0xa320} at compile time, so the argument that bytecode is unchanged
is a compile-time substitution argument, confirmed by the pinned hash, not an
empirical one.

\section{Fix 2 --- Semantic local names over \texttt{varN}}
\label{sec:read-varnames}
\vstatus{OPEN}

\paragraph{The problem.} 163 distinct local names \texttt{var0}\ldots\texttt{var162}
(\texttt{next\_var}, \texttt{yul\_\allowbreak{}emit.rs:1125--1129}), declared 442 times
because the counter restarts per emitted block, so the same name denotes
different values in different blocks. A reader cannot tell
\texttt{var47} from \texttt{var48}.

\paragraph{The fix is already half-built.} The generator already has
\texttt{column\_eval\_var} (\texttt{yul\_emit.rs:1142}), which produces
semantic names like \texttt{f\_3}, \texttt{a\_7\_next\_1}. The infrastructure
for meaningful names exists; it is just not used for the generic
\texttt{init\_var} path.

\paragraph{Evidence the hint mechanism already works.} In
\texttt{yul\_emit.rs}, \texttt{init\_var} (L1113) takes an
\texttt{Option<String>} name and 23 call sites exercise it: \textbf{4 pass a
real hint} (\texttt{Some(var\_name)} at L1041/L1053 for fixed/advice evals,
\texttt{Some(column\_\allowbreak{}eval\_\allowbreak{}var("i", col, rot))} at L1059, and
\texttt{Some(format!("c\_\allowbreak{}\{\}", \ldots))} for challenges at L1062--1065),
while \textbf{19 pass \texttt{None}} and fall back to \texttt{varN} (L727,
904, 917, 927, 966, 973, 976, 989, 993, 1003, 1005, 1014, 1025, 1033, 1068,
1073, 1079, 1085, 1088). The hinted sites are exactly the leaf reads; the
\texttt{None} sites are the generic arithmetic ---
\texttt{mulmod}/\texttt{addmod}/negation/\texttt{q\_pow5}/constant
materialization. So the mechanism is proven in production; the readability
gap is precisely those 19 generic-arithmetic sites.

\paragraph{Concrete change.} Thread a \emph{hint} through
\texttt{init\_var} and fall back to \texttt{varN} only when no hint is
available. The useful hints are not protocol-level names:
\texttt{init\_var} only emits \emph{gate-expression subterms} (the
vanishing, Lagrange and pairing values live in templates and
\texttt{kzg/mod.rs} and are already named, e.g.\ \texttt{x\_n},
\texttt{l\_0}, \texttt{x4\_pow\_5}). The hint should come from the gate and
constraint names the generator already carries in the quotient identity
manifest (\texttt{builder/\allowbreak{}api.rs:121}), plus the term's role, and must be
uniquified, because Yul rejects a second \texttt{let} of a name visible in
the same scope:

\begin{lstlisting}[style=rust]
// yul_emit.rs: init_var already takes Option<String> for the name.
// Callers that know the gate/constraint should pass a unique hint:
let hint = format!("g{gate}_c{constraint}_prod{}", self.next_suffix());
let (lines, prod) = self.init_var(prod_expr, Some(hint));
// next_var() stays as the last-resort fallback.
\end{lstlisting}

\paragraph{Cost / risk.} Low. \texttt{init\_var} already accepts an
optional name; this is a call-site change. The var-cache keys on the
\emph{value}, not the name (L1113--1122), so CSE behavior is unchanged.
Uniqueness of hints is the one new obligation.

\section{Fix 3 --- Decompose \texttt{verifyProof} into named functions}
\label{sec:read-decompose}
\vstatus{OPEN}

\paragraph{The problem.} \texttt{verifyProof} is 3{,}286 lines of inline
Yul. No reviewer holds that in working memory. The protocol has
$\sim$12 named stages (parse, challenges, instance-eval, identity
families, linearization, multi-open, final pairing, accumulator) but they
are not visible as units.

\paragraph{Concrete change.} The renderer already emits per-stage blocks
(Ch.~\ref{ch:solidity} documents ``PCS blocks'', ``Final pairing'',
etc.). Wrap each stage in a Yul \texttt{function} with the protocol name.
Stages should communicate through the existing memory map
(\texttt{THETA\_MPTR}\ldots\texttt{X4\_MPTR}, \texttt{FINAL\_COM\_MPTR},
\ldots) and return only a success flag. Returning nine challenges from one
function would add exactly the stack pressure the split is meant to remove.

\begin{lstlisting}[style=yul]
function verify_parse_proof(proof_cptr) -> ok { ... }
function verify_challenges() -> ok { ... }        // writes THETA_MPTR..X4_MPTR
function verify_instance_eval() -> ok { ... }     // writes INSTANCE_EVAL_MPTR
function verify_identity_families() -> ok { ... } // writes QUOTIENT_EVAL_MPTR
function verify_kzg_multiopen() -> ok { ... }     // writes FINAL_COM_MPTR, V_MPTR
function verify_final_pairing() -> ok { ... }
\end{lstlisting}

The top-level \texttt{verifyProof} becomes a 15-line table of calls that
\emph{reads like the protocol spec}. Each function is independently
reviewable and independently testable against the Rust oracle. It also
bounds each function's CFG, which is what stock slither's SSA pass needs
(\S\ref{sec:c5-tools}).

\paragraph{The seam already exists in the assembly pipeline.} The contract
body is not one monolithic string --- it is assembled from a structured
\texttt{QuotientComputationBlocks} collection
(\texttt{plan.rs:68} \texttt{into\_\allowbreak{}template\_\allowbreak{}parts()}) plus
\texttt{kzg::computations()} (\texttt{artifacts.rs:110--117}). The
quotient stages are already separated into a typed block collection before
they are concatenated into the single \texttt{assembly ("memory-safe")}
body (the test \texttt{generated\_\allowbreak{}heavy\_\allowbreak{}assemblies\_\allowbreak{}keep\_\allowbreak{}memory\_\allowbreak{}safe\_\allowbreak{}for\_\allowbreak{}via\_\allowbreak{}ir},
\texttt{tests.rs:2127}, asserts that annotation). So
Fix~3 is a mechanical transformation at the assembly seam: wrap each
element of the block collection in a named \texttt{function} and replace
the inline concatenation with a call sequence. The decomposition does not
require re-deriving the stages --- they are already discrete.

\paragraph{Cost / risk.} Medium. Yul function calls add a small stack
frame; via-IR's inliner will probably re-inline single-call functions, so gas
should be close to neutral, but that is not guaranteed and the bytecode will
change. Two further caveats. First, the restructure changes which locals
solc spills into $[\texttt{0x80},\textit{memoryguard})$, an area the
generated code also writes (\S\ref{sec:c11-build}), so re-check the spill
area and live ranges after the change. Second, it does not by itself remove
the via-IR requirement: the first legacy-codegen stack-too-deep error is in
the helper \texttt{load\_\allowbreak{}acc\_\allowbreak{}coord\_\allowbreak{}shifted} (L573). This is the
largest change but the single biggest auditability win.

\section{Fix 4 --- Annotate the fused \texttt{q\_com} MSM pairs}
\label{sec:read-qcom}
\vstatus{PARTIAL}

From \S\ref{sec:pr-qcom}: the quotient commitment is fused into the
78-pair final MSM (63 multi-open terms, 14 linearization terms --- 4 quotient
limbs and 10 simple selectors --- and \texttt{f\_com}), and
\texttt{Q\_COM\_MPTR} is declared-but-unused. Two cheap fixes:

\begin{enumerate}
\item \textbf{Drop the dead constant.} \texttt{Q\_COM\_MPTR}
(\texttt{L137}) is never written --- confirmed: \textbf{0}
\texttt{mstore}/\texttt{mload}/\texttt{mcopy} operands reference it.
\vstatus{ADDRESSED} for the comment half: L134--136 already say it
``intentionally aliases \texttt{Q\_\allowbreak{}EVAL\_\allowbreak{}SET\_\allowbreak{}MPTR} and has zero reserved
capacity until a future emitter starts writing \texttt{Q\_COM\_MPTR}''.
Deleting it (as the unmerged \texttt{41b78eb9} does) remains optional. The
same applies to the four other dead pointer constants
(\texttt{CHALLENGE\_MPTR}, \texttt{G1\_\allowbreak{}IDENTITY\_\allowbreak{}MPTR},
\texttt{QUOTIENT\_\allowbreak{}RETURN\_\allowbreak{}MPTR}, \texttt{TRACE\_U256\_MPTR};
\S\ref{sec:c5-tools}). \texttt{G1\_\allowbreak{}IDENTITY\_\allowbreak{}MPTR}'s comment (L144--148)
describes ``the four \texttt{mload}s below'', which this render does not
emit.
\item \textbf{Annotate the MSM pairs.} \vstatus{OPEN}: the staging block
(L3298--3458) has no per-pair comments. In the final-MSM input build,
emit a comment per pair: \texttt{/\allowbreak{}/\allowbreak{} pair 12: q\_\allowbreak{}com contribution
(x4\_\allowbreak{}pow\_\allowbreak{}4)}. With Fix~1's named addresses this becomes self-documenting.
\end{enumerate}

\section{Fix 5 --- Emit a generated address-map appendix}
\label{sec:read-addressmap}
\vstatus{OPEN}

Even with Fix~1, some addresses remain raw hex. The generator should emit
a machine-generated legend as a comment block at the top of the contract
(or a companion \texttt{.md} file):

\begin{lstlisting}[style=plain]
// MEMORY MAP (generated; do not hand-edit)
//   0x0064  PROOF_CPTR          proof calldata base
//   0x0080  TRANSCRIPT_MPTR     transcript buffer (phase: transcript;
//           = RETURN_MPTR, QUOTIENT_RETURN_MPTR in later phases)
//   0x0080..0x08e0              solc via-IR spill area (memoryguard)
//   0x1ec4  NUM_INSTANCE_CPTR   instance count word
//   0x1ee4  INSTANCE_CPTR       instance array base
//   0x4140  q_program           VM bytecode blob (0x11cf bytes)
//   0x6ea0  FINAL_COM_MPTR      KZG reduced commitment
//   0xa320  PCS_FINAL_MSM_MPTR  final MSM input scratch (78 pairs;
//           quotient-VM scratch in an earlier phase)
//   ...
\end{lstlisting}

This is free to generate (the \texttt{MemoryMap} already has every
\texttt{(name, start, len, lifetime)}, \texttt{memory.rs:178--186}; the
production contract already declares 68 named pointer constants that could
seed the legend) and gives a reviewer a single lookup table. It turns
``what is 0xa320?'' from a grep hunt into a one-line answer. Two details
matter: show each region's lifetime, because addresses are reused across
phases, and include the compiler's spill area from the compiled
\texttt{memoryguard}, which would make the C11 overlap
(\S\ref{sec:c11-build}) visible at a glance.

\section{Fix 6 --- Collapse the five near-identical Horner blocks}
\label{sec:read-horner}
\vstatus{OPEN}

The \texttt{f\_eval} section (\S\ref{sec:pr-feval}) emits five
reversed point-set blocks. In the production contract the
\texttt{f\_eval} block spans L3107--3269 and contains five
\texttt{f\_\allowbreak{}eval := addmod(mulmod(f\_\allowbreak{}eval, x2, r), eval, r)} steps
(L3158, L3202, L3231, L3260, L3267). The Horner step itself is one line;
what repeats is the per-set Lagrange interpolation that precedes it. Its
size depends on the set's cardinality (3, 3, 2, 2, 1 at L3116, L3160, L3204,
L3233, L3262), so a \texttt{horner(coeffs, x)} helper would remove almost
nothing. The generator should instead emit one loop-based helper,
parameterized by the set's point count and its \texttt{ROT\_POINTS}/\texttt{Q\_EVAL\_SET}
offsets, and call it five times. This removes $\sim$150 lines of
near-copy-paste and makes the ``five sets'' structure explicit rather than
something the reviewer must pattern-match. It is not bytecode-neutral: a
loop costs a little more than unrolled code. It also exposes a gas win:
each block calls \texttt{scalar\_inv} (a modexp) once, and one batched
inversion across the five sets would save four modexps ($\approx$5.4k gas under
EIP-2565, $\approx$16.3k under EIP-7883).

\section{Prioritized summary}
\label{sec:read-priority}

\begin{center}
\begin{tabular}{@{}cp{7.0cm}llp{2.3cm}@{}}
\toprule
\textbf{\#} & \textbf{Fix} & \textbf{Impact} & \textbf{Effort} & \textbf{Status} \\
\midrule
1 & Named memory addresses (symbolic \texttt{Value::Identifier} at the emit site) & High & Low & \vstatus{OPEN} \\
2 & Semantic local names (reuse \texttt{init\_var} hint) & High & Low & \vstatus{OPEN} \\
3 & Decompose \texttt{verifyProof} into named functions & High & Med & \vstatus{OPEN} \\
4 & Annotate fused \texttt{q\_com} MSM pairs & Med & Low & \vstatus{PARTIAL} \\
5 & Generated address-map appendix & Med & Trivial & \vstatus{OPEN} \\
6 & Loop helper for the five per-set interpolation blocks & Med & Low & \vstatus{OPEN} \\
\bottomrule
\end{tabular}
\end{center}

\paragraph{The unifying principle.} Fixes 1, 2, and 5 all share one
property: \emph{the information is already in the generator; the renderer
just discards it.} A single design change --- ``every emitted address and
local carries its semantic name, with hex as fallback'' --- recovers most
of the lost readability without touching a line of protocol logic. For
these three, the compiled bytecode is expected to be unchanged and only the
source text gains names. That is why they are low-risk: they cannot
introduce a soundness bug, because they cannot change what the contract
\emph{does} --- only what it \emph{looks like} to the human reading it. The
pinned runtime hash makes that claim testable. Fixes 3 and 6 do change the
code's structure: they need the full differential suite and the spill-area
re-check of \S\ref{sec:c11-build}.

\paragraph{What this does not fix.} The fuzz-coverage gap
(\S\ref{sec:pr-fuzz-coverage}) and the hybrid VM+Yul model
(\S\ref{sec:pr-vm-hybrid}) are not readability problems and are not
addressed here. Readability makes the contract \emph{auditable}; it does
not make the production shape \emph{tested}. Both matter, and they are
orthogonal workstreams.
